\scanpage{110}
\mspage{109}
\[
S\mathcal{T}_G\overset{\mathrm{déf}}=\unsure{\Pi_1|\pi_1}(SP_G)\xleftarrow{\sim}\unsure{\Pi_1|\pi_1}(\widetilde{P}_G^B).
\]

D'autre part, on a un diagramme cartésien évident
\diagram{\begin{tikzcd}[row sep=large,column sep=large]
\widetilde{P}_G^B \arrow[r] \arrow[d,"{\begin{gathered}\mathbb U^B\text{-tors.}\\\text{sur }\widetilde M_G^C\end{gathered}}"'] & SP_{G'} \arrow[d] \\
\widetilde{M}_G^C \arrow[r] & M_{G'}
\end{tikzcd}}
qui exprime simplement que pour $X'$ construit à l'aide de $(X,(D_{\widetilde a}))$, le \deleted{\unsure{t}} $\mathbb U^B$-torseur $P_{B,X}$ \deleted{$\unsure{T_G}$ est} est can. isom. (on dirait même, \uline{égal}) au $\mathbb U^B$-torseur $P$. On trouve ainsi un morphisme
\diagram{\begin{tikzcd}[column sep=large,row sep=small]
\unsure{\Pi_1|\pi_1}(\widetilde{P}_G^B) \arrow[r] \arrow[d,equal] & \unsure{\Pi_1|\pi_1}(SP_{G'}) \arrow[d,equal] \\
S\mathcal{T}_G & S\mathcal{T}_{G'}
\end{tikzcd}}

Considérons le monoïde topologique multiplicatif $]0,1]$ (avec élément unité $1$) et sa puissance $C$-ième
\[
\unsure{\mathfrak m|m}_{\vec C}=]0,1]^{\vec C},
\]
\margin{il vaut mieux prendre $]0,1]^C$, sinon il y a des ``paramètres bidon''}
alors $m_C$ opère sur le $M_{\unsure{e|G}}$-fibré $\widetilde{M}_G^C$ de façon évidente (homothétie sur les disques $D_{\widetilde a}$),
\margin{on prend $\rho_{\widetilde a'}=\rho_{\widetilde a''}$ pour $\widetilde a',\widetilde a''$ sur $a\in C$}
